\chapter*{Bảng ký hiệu và chữ viết tắt}
\addcontentsline{toc}{chapter}{\bfseries Bảng ký hiệu và chữ viết tắt}

\begin{tabular}
    {@{\hspace{0.6cm}} l @{\hspace{1.8cm}}p{11.5cm}l}
    $\mathcal{H}$ & không gian Hilbert thực\\
    $\mathbb R$ & tập các số thực\\
    $\mathbb R^n$ & không gian Euclid $n$-chiều\\
    $l^2$ & không gian các dãy bình phương khả tổng\\
    $L^2[0,1]$ & không gian các hàm bình phương khả tích trên đoạn [0,1]\\
    $\langle x,y\rangle$ & tích vô hướng của $x$ và $y$\\
    $\mathbf{0}$ & phần tử không\\
    $\|x\|$ & chuẩn của $x$\\
    $\varnothing$ & tập rỗng\\
    $\mathcal{C}^\complement$ & bao lồi của tập hợp C\\
    $\overline{\mathcal{C}}$ & phần bù của tập hợp C\\
    $\|\mathcal{F}\|$ & chuẩn của toán tử $\mathcal{F}$\\
    $\mathcal{F}^*$ & toán tử liên hợp của toán tử $\mathcal{F}$\\
    Fix$(\mathcal{F})$ & tập điểm bất động của toán tử $\mathcal{F}$\\
    $I^{\mathcal{H}}$ & toán tử đồng nhất của không gian $\mathcal{H}$\\
    $P_\mathcal{C}(x)$ & hình chiếu mêtric của phần tử $x$ lên tập $\mathcal{C}$\\
    $x^n \to x^*$ & dãy $\{x^n\}$ hội tụ mạnh đến $x^*$\\
    $x^n \rightharpoonup x^*$ & dãy $\{x^n\}$ hội tụ yếu đến $x^*$\\
    VIP($\mathcal{G}, \Omega$) & bài toán bất đẳng thức biến phân với ánh xạ giá $\mathcal{G}$ và tập ràng buộc $\Omega$\\
    SFP & bài toán chấp nhận tách\\
    SVIP & bài toán bất đẳng thức biến phân tách
\end{tabular}
